\documentclass[conference]{IEEEtran}
\IEEEoverridecommandlockouts

\usepackage[brazil]{babel}
\usepackage[utf8]{inputenc}
\usepackage[T1]{fontenc}
\usepackage{cite}
\usepackage{amsmath,amssymb,amsfonts}
\usepackage{algorithmic}
\usepackage{graphicx}
\usepackage{textcomp}
\usepackage{xcolor}
\usepackage{booktabs}
\usepackage{url}

\def\BibTeX{{\rm B\kern-.05em{\sc i\kern-.025em b}\kern-.08em
    T\kern-.1667em\lower.7ex\hbox{E}\kern-.125emX}}

\begin{document}

\title{Caracterização de Padrões de Depreciação Automotiva sob a Transição Energética: Uma Abordagem Baseada em Transformação Bidimensional de Séries Temporais}

\author{\IEEEauthorblockN{Gabriel Martins Brum}
\IEEEauthorblockA{\textit{Departamento de Ciência da Computação e Estatística} \\
\textit{Universidade Estadual Paulista (UNESP)}\\
São José do Rio Preto, Brasil \\
g.brum@unesp.br}
}

\maketitle

% o abstract deve ter uma estrutura assim. contexto -> problema -> o que o trabalho fez pra resolver o problema -> resultados encontrados
\begin{abstract}
A transição para a mobilidade elétrica introduz incertezas na precificação de veículos seminovos em razão da perda de capacidade de baterias, rápida obsolescência técnica e guerras comerciais de preços. Este artigo investiga a desvalorização veicular a partir de microdados da base fipeX. A investigação analisa se veículos elétricos apresentam desvalorização acelerada em relação a modelos convencionais de mesma faixa de preço de lançamento, bem como se a codificação de séries temporais em matrizes bidimensionais via \textit{Gramian Angular Fields} e \textit{Markov Transition Fields} aprimora o agrupamento não supervisionado frente a abordagens clássicas. A metodologia emprega a Taxa de Retenção de Valor em uma janela temporal de 32 meses com 4.457 séries, redução de dimensionalidade por PCA e modelagem por K-Means e DBSCAN. A validação das partições foi realizada por métricas intrínsecas e análise de variância bidirecional.
\end{abstract}

\begin{IEEEkeywords}
Mineração de Séries Temporais, Veículos Elétricos, \textit{Reshaping}, Agrupamento, Depreciação Automotiva.
\end{IEEEkeywords}

\section{Introdução}

% contextualização deve ter dois eixos, um do mercado automotivo e outro do reshaping de séries temporais para imagem para aumento de desempenho de modelos classificadores. o contexto do mercado automotivo deve apresentar as pesquisas que falem sobre os estudos que investigam a depreciação de veiculos elétriso/híbridos e comparam com veículos a combustão (porem nenhum eh realizado no brasil). ja o contexto do reshaping de series temporais devem apresentar estudos e casos que a transformacao para o dominio das imagens trouxe beneficios, e isso eh uma tecnica recente na literatura.

% motivação e justificativa: dado que o mercado brasileiro dos veículos elétricos e híbridos é recente, e carece de estudos que investiguem a depreciação dos mesmos, a literatura apresenta uma lacuna. o reshaping serve tambem para uma melhor caracterizacao dos padroes de depreciacao. alem de que o reshaping de series temporais de veiculos tambem seja uma lacuna (checar se é mesmo).

% objetivos: alem de apenas caracterizar a depreciacao, o trabalho busca investigar se o mapeamento para o dominoi das imagens melhora a separabilidade dos grupos, e se o agrupamento une veiculos mais por faixa de preço ou por matriz energetica. tambem analisarei se a depreciacao dos carros a combustao foram afetados pelos eletricos ou nao.

% contribuicoes: (1) caracterizacao da depreciacao de veiculos eletricos e hibridos no brasil, (2) investigacao do efeito da matriz energetica sobre a depreciacao frente a faixa de preco, (3) investigacao do efeito do reshaping sobre a separabilidade dos grupos. (4) investigacao do efeito da entrada de veiculos eletricos sobre a depreciacao dos veiculos a combustao.

O mercado automotivo global e brasileiro atravessa um processo de reorganização tecnológica impulsionado pela introdução e expansão de veículos elétricos (EVs) e híbridos \cite{alvarez2024modelagem}. A especulação em torno da depreciação desses veículos cria uma barreira para potenciais compradores, gerando receio quanto à revenda e elevando a incerteza sobre a taxa de desvalorização dos ativos \cite{rush2022vehicle}. Automóveis equipados com motores a combustão apresentam trajetória de desvalorização consolidada e linearizada em função da idade do bem \cite{storchmann2004depreciation}. Em contrapartida, os veículos eletrificados enfrentam incertezas operacionais singulares, destacando-se a perda de capacidade de armazenamento das baterias de íons de lítio e a rápida substituição técnica por modelos de maior autonomia \cite{gautam2024depreciation}. A avaliação empírica conduzida sobre cinco mercados desenvolvidos apurou depreciação média anual de 13,9\% para veículos elétricos (1,16\% ao mês), superando a taxa de 10,4\% ao ano observada em modelos convencionais a gasolina \cite{schloter2022empirical}. A literatura sobre o comportamento de depreciação veicular e políticas tarifárias tem se concentrado prioritariamente em economias desenvolvidas e na China, subsistindo uma lacuna empírica quanto à dinâmica do mercado brasileiro sob a transição energética \cite{li2026global, silva2024avaliacao}.

Para além da dinâmica setorial de mercado, a análise e a descoberta de padrões comportamentais em séries temporais têm incorporado técnicas de codificação matricial bidimensional (\textit{reshaping}) \cite{tan2019introduction, wang2015imaging}. Abordagens convencionais representam o histórico de preços como vetores unidimensionais em $\mathbb{R}^T$, empregando distâncias Euclidianas sensíveis a desvios de escala e ruídos pontuais. O algoritmo Dynamic Time Warping (DTW) flexibiliza o alinhamento de sequências desfasadas, mas distensões no eixo temporal tendem a suprimir a sincronia de choques macroeconômicos exógenos, além de introduzir elevado custo computacional quadrático sobre bases volumosas \cite{sakoe1978dynamic}. Em contraposição, a conversão de séries temporais em matrizes bidimensionais atua como um agregador de atributos intrínsecos sobre a sequência original sem demandar variáveis externas. As transformações Gramian Angular Summation Field (GASF), Gramian Angular Difference Field (GADF) e Markov Transition Field (MTF) proporcionam reduções de 12\% a 48\% no erro de classificação frente a vetores brutos ao preservar relações angulares e matrizes de transição, promovendo adicionalmente maior estabilidade e ganho de silhueta no agrupamento financeiro \cite{wang2015imaging, wu2023imaging}. Esse panorama delimita duas lacunas centrais investigadas neste artigo: a carência de estudos empíricos sobre a Taxa de Retenção de Valor (VRR) de veículos elétricos e híbridos em comparação a modelos a combustão no mercado brasileiro e a inexistência de investigações que empreguem técnicas de \textit{reshaping} para a descoberta de padrões de desvalorização veicular.

\subsection{Objetivos, Questões de Pesquisa e Contribuições}

O objetivo deste estudo consiste em analisar a dinâmica temporal de desvalorização veicular sob a transição energética no Brasil por meio de agrupamento não supervisionado, contrastando espaços vetoriais unidimensionais com representações matriciais bidimensionais combinadas à redução dimensional. A investigação é guiada por duas questões centrais de pesquisa:
(1) a entrada de veículos elétricos induz desvalorização relativa mais acelerada em comparação a veículos convencionais pareados no mesmo estrato econômico de lançamento, ou o padrão observado reflete prioritariamente a faixa de preço do ativo?
(2) a codificação de séries temporais em matrizes bidimensionais projetadas via PCA produz agrupamentos com maior coesão e separabilidade intrínseca em relação às métricas vetoriais clássicas Euclidianas e DTW?

A partir dessa formulação, o trabalho apresenta três contribuições principais:
(1) a caracterização quantitativa da retenção de valor de automóveis elétricos, híbridos e a combustão no mercado brasileiro ao longo de 32 meses de acompanhamento com microdados da base fipeX;
(2) a comprovação estatística do papel preponderante da matriz energética sobre a depreciação por meio de pareamento econômico intra-estrato e análise de variância bidirecional;
(3) a validação da codificação matricial GASF associada ao PCA como alternativa computacionalmente eficiente e com qualidade de separabilidade superior ao alinhamento elástico clássico.

\section{Materiais e Métodos}

O procedimento metodológico foi dividido em quatro etapas executadas sequencialmente: preparação da base de dados, pré-processamento das séries temporais de retenção de valor, codificação matricial bidimensional com redução dimensional e modelagem de agrupamento com validação das partições.

\subsection{Base de Dados FipeX}

Os dados foram obtidos do repositório público fipeX, que reúne o histórico de preços médios mensais de veículos automotores comercializados no Brasil apurados pela Fundação Instituto de Pesquisas Econômicas (FIPE) entre maio de 2005 e janeiro de 2026 \cite{fipex2024}. A base original continha 9.276.121 registros cadastrais com atributos de identificação do modelo, montadora, ano de modelo, tipo de veículo, matriz energética e preço médio de mercado cotado em Reais ($P_{i,t}$).

\subsection{Pré-Processamento e Análise Exploratória}

O procedimento de limpeza e preparação foi conduzido para eliminar inconsistências e isolar a evolução dos preços em relação à idade cronológica dos veículos. Foram descartadas 83 instâncias duplicadas e 22 registros com valor monetário nulo. A base foi restrita a automóveis de passeio, o que resultou na exclusão de 3.565.306 registros referentes a motocicletas, caminhões e utilitários de grande porte. Modelos antigos cadastrados com atraso na FIPE foram eliminados por meio da restrição temporal $(ano\_lancamento - ano\_modelo) \le 1$, descartando 1.942.304 registros para assegurar que a idade zero corresponda ao período de estreia comercial do ativo. Foram removidos 52.721 registros de combustíveis com volume residual (etanol e gás natural) e 24.996 registros de montadoras com baixa participação de mercado. Esse processo totalizou a exclusão de 5.585.432 registros, correspondendo a uma redução de 60,2\% da base inicial.

A base remanescente de 3.690.689 registros pertencentes a 63 montadoras foi submetida ao cálculo da Taxa de Retenção de Valor (VRR), definida pela relação entre o preço no mês de avaliação $t$ e o preço de lançamento no mês inicial:
\begin{equation}
VRR_{i,t} = \frac{P_{i,t}}{P_{i,0}}, \quad t \in \{0, 1, \dots, T-1\}.
\label{eq:vrr}
\end{equation}
O cálculo da Equação~\ref{eq:vrr} remove a magnitude nominal dos preços em Reais e permite a comparação direta entre veículos de diferentes faixas de mercado. Para viabilizar a modelagem, a amostra foi restrita a veículos lançados a partir de 2019 com acompanhamento contínuo por uma janela de $T = 32$ meses (2,6 anos), tolerando até 10\% de interpolação linear para meses ausentes. Esse corte consolidou uma matriz balanceada de 4.457 séries temporais completas. A amostra resultante é composta por 1.756 veículos Flex (39,4\%), 1.258 a Gasolina (28,2\%), 871 a Diesel (19,5\%), 382 Híbridos (8,6\%) e 190 Elétricos (4,3\%).

A categorização econômica foi estabelecida a partir dos quartis gerais dos preços de lançamento, apurando $Q_1$ igual a R\$ 82.998,00, mediana de R\$ 155.043,00 e $Q_3$ de R\$ 325.283,00. Os veículos foram classificados nos estratos de Entrada ($P_{i,0} \le Q_1$), Intermediário ($Q_1 < P_{i,0} \le Q_3$) e Luxo ($P_{i,0} > Q_3$).

\subsection{Codificação Bidimensional de Séries Temporais}

As trajetórias temporais de retenção de valor foram convertidas em matrizes quadradas de dimensão $32 \times 32$ sem redimensionamento prévio, conforme o método proposto por Wang e Oates \cite{wang2015imaging}. A série normalizada no intervalo $[-1, 1]$ foi inicialmente projetada em coordenadas polares com ângulo $\phi_t = \arccos(\widetilde{VRR}_t)$ e raio $r_t = t/T$. A matriz Gramian Angular Summation Field (GASF) foi gerada a partir da soma cosseno entre os pontos temporais:
\begin{equation}
G_{(i,j)} = \cos(\phi_i + \phi_j) = \widetilde{x}_i \widetilde{x}_j - \sqrt{1 - \widetilde{x}_i^2}\sqrt{1 - \widetilde{x}_j^2}.
\label{eq:gasf}
\end{equation}
A transformação descrita na Equação~\ref{eq:gasf} preserva a correlação temporal e os valores médios de retenção do ativo. A matriz Gramian Angular Difference Field (GADF) foi calculada por meio da diferença seno:
\begin{equation}
D_{(i,j)} = \sin(\phi_i - \phi_j) = \widetilde{x}_j \sqrt{1 - \widetilde{x}_i^2} - \widetilde{x}_i \sqrt{1 - \widetilde{x}_j^2},
\label{eq:gadf}
\end{equation}
agregando informação relativa às taxas locais de variação e à velocidade de perda de valor ao longo da série. A técnica Markov Transition Field (MTF) foi construída discretizando as trajetórias em cinco quantis de amplitude e computando as probabilidades de transição estocástica de primeira ordem entre os passos temporais $i$ e $j$. Como o vetor achatado de cada imagem contém 1.024 atributos por canal e 3.072 atributos na fusão dos três canais, a técnica de Análise de Componentes Principais (PCA) foi aplicada retendo 95\% da variância explicada acumulada \cite{tan2019introduction}. Essa projeção reduziu os dados para 2 componentes no GASF e 5 componentes no GADF, eliminando ruídos colineares.

\subsection{Modelagem de Agrupamento e Protocolo de Validação}

O algoritmo K-Means foi avaliado com inicialização K-Means++ variando $K$ de 2 a 7 sobre o espaço 1D VRR Euclidiano, sobre o espaço 1D com alinhamento temporal e sobre os espaços latentes gerados pelas imagens 2D via PCA \cite{arthur2007k, sakoe1978dynamic}. O algoritmo baseado em densidade DBSCAN também foi investigado, adotando o raio $\varepsilon$ determinado pelo ponto de inflexão das curvas de $k$-distâncias para $MinPts = 5$ \cite{ester1996density, tan2019introduction}.

A avaliação das partições foi conduzida por critérios intrínsecos e testes extrínsecos. Na validação intrínseca, foram calculados o coeficiente de Silhueta médio e o índice Davies-Bouldin, acompanhados da análise do decaimento da soma dos erros quadráticos \cite{rousseeuw1987silhouettes, davies1979cluster}. Na validação extrínseca com os rótulos de combustível e faixa de preço, foram computados o índice Rand ajustado (ARI), a informação mútua normalizada (NMI), a homogeneidade, a completude e o coeficiente $V$ de Cramér. A contribuição de cada fator sobre a variação líquida aos 32 meses ($\Delta VRR = VRR_{32m} - VRR_{0m}$) foi quantificada por análise de variância (ANOVA) bidirecional ($2 \times 5$).

\section{Resultados e Discussão}

\subsection{Análise Exploratória e Pareamento de Segmentos}

A amostra de 4.457 trajetórias de retenção expõe o comportamento singular do mercado brasileiro entre 2019 e 2026. Em contraste com o padrão histórico de depreciação linear, automóveis a combustão registraram valorização nominal atípica motivada por choques na cadeia global de suprimentos, atingindo retenção média superior a 115\% em modelos Flex e Diesel aos 12 e 24 meses. Em contrapartida, veículos elétricos exibiram desvalorização contínua. Como 93,2\% dos elétricos pertencem aos estratos Intermediário e de Luxo (mediana de R\$ 322.000,00) e 75,5\% dos modelos Flex integram o estrato de Entrada (mediana de R\$ 77.000,00), a avaliação comparativa demanda pareamento intra-estrato para isolar a tecnologia de propulsão da magnitude nominal do preço inicial.

\subsection{Eficácia Comparativa das Representações 1D e 2D}

A qualidade do agrupamento pelo algoritmo K-Means foi avaliada para partições de $K=2$ a $K=7$ sobre os espaços vetoriais 1D e os espaços latentes das transformações bidimensionais reduzidas por PCA. Para a partição com quatro grupos ($K=4$), a projeção matricial GASF alcançou a maior coesão e separabilidade intrínseca ($s = 0,410$, $DB = 0,784$, tempo de 0,05 s), superando expressivamente o vetor 1D Euclidiano ($s = 0,311$, $DB = 1,028$, 0,06 s) e o alinhamento DTW com DBA ($s = 0,243$, $DB = 1,212$, 71,62 s). A transformação GADF produziu separabilidade semelhante ($s = 0,424$, $DB = 0,967$, 0,05 s), enquanto a codificação MTF obteve desempenho insatisfatório ($s = 0,064$, $DB = 3,116$, 0,80 s) pela supressão do gradiente contínuo na discretização em quantis. A fusão multicanal registrou $s = 0,200$, $DB = 1,502$ e 0,24 s, denotando dispersão de atributos. O processamento sobre os dois componentes do GASF consumiu apenas 0,05 segundos, estabelecendo uma aceleração superior a 1.400 vezes frente ao cálculo elástico do DTW.

\subsection{Arquétipos de Depreciação e Validação Extrínseca}

O particionamento no espaço GASF com $K=4$ identificou quatro arquétipos temporais de retenção: Valorização Moderada ($C_1$, média de $+23,3\%$ aos 32 meses, 1.005 veículos), Depreciação Acelerada ($C_2$, média de $-13,6\%$, 1.624 veículos), Super-Valorização ($C_3$, média de $+39,6\%$, 338 veículos) e Estabilidade ($C_4$, média de $+5,7\%$, 1.490 veículos). O arquétipo $C_2$ reuniu 65,8\% de todos os veículos elétricos da base (125 unidades) e 48,4\% dos modelos híbridos, ao passo que $C_3$ não conteve nenhum modelo elétrico.

Os testes de validação extrínseca registraram $ARI = 0,012$, $NMI = 0,022$ e $V = 0,140$ para combustível, e $ARI = 0,012$, $NMI = 0,024$ e $V = 0,159$ para faixa de preço. Esses escores próximos de zero confirmam que o agrupamento não supervisionado não reproduz categorias cadastrais prévias, mas organiza os veículos pela dinâmica temporal de suas trajetórias.

A análise de variância bidirecional ($2 \times 5$) sobre a variação líquida aos 32 meses indicou que o tipo de combustível responde por quase o dobro da variabilidade explicada frente ao estrato econômico ($SS = 142.704$ contra $SS = 73.614$, ambos com $p < 10^{-39}$). No estrato de Luxo, com veículos pareados no mesmo patamar de entrada, modelos elétricos registraram perda mediana de $-23,8\%$, enquanto modelos a gasolina obtiveram ganho de $+1,2\%$ e a diesel alcançaram $+12,1\%$. Essa diferença de 25 pontos percentuais confirma que a taxa de desvalorização decorre da tecnologia de propulsão e não do patamar de preço.

\section{Conclusão}

Este estudo investigou a retenção de valor veicular durante a transição energética no Brasil por mineração de séries temporais. Em resposta à primeira questão de pesquisa, a desvalorização acelerada em veículos elétricos decorre da tecnologia de propulsão e não apenas do estrato de preço. Mais de 65\% dos modelos elétricos concentram-se no grupo de desvalorização contínua, e o teste ANOVA comprovou a primazia do combustível sobre o estrato de mercado ($SS = 142.704$ contra $SS = 73.614$), registrando defasagem de 25 pontos percentuais frente a veículos a gasolina de mesmo valor inicial.

Em resposta à segunda questão de pesquisa, a transformação bidimensional GASF com PCA reduziu a dispersão dos dados e elevou o coeficiente de Silhueta de 0,311 para 0,410, viabilizando o agrupamento com rapidez computacional três ordens de grandeza superior ao método DTW. Como limitação metodológica, o corte temporal de 32 meses exclui modelos lançados muito recentemente. Trabalhos futuros devem estender o monitoramento para avaliar a consolidação das redes de reciclagem e substituição de baterias.

\bibliographystyle{IEEEtran}
\bibliography{references}

\end{document}
